\section{Exact three-dimensional Matignon threshold}
\label{sec:exact-3x3}

Let \(-A\) be strict P and use the notation of Lemma~\ref{lem:simplex-reduction}. For \(D\succ0\),
\[
\det(\lambda I-DA)
=
\lambda^3+a_D\lambda^2+b_D\lambda+c_D,
\]
with
\[
a_D=s,\qquad
b_D=s^2B_\beta(x),\qquad
c_D=s^3\kappa x_1x_2x_3.
\]

\subsection{Exact cubic boundary}

\begin{lemma}[Exact cubic Matignon boundary]
\label{lem:cubic-boundary}
Let
\[
p(\lambda)=\lambda^3+a\lambda^2+b\lambda+c,
\qquad a,b,c>0,
\]
and \(\pi/3<\theta<\pi/2\). Set
\[
u=\cos\theta,\qquad K=1-4u^2>0,
\]
\[
r_\theta(a,b)
=
\frac{ua+\sqrt{u^2a^2+Kb}}{K},
\]
and
\[
H_\theta(a,b)
=
r_\theta(a,b)^2
\left(a+2u\,r_\theta(a,b)\right).
\]
Then all roots satisfy
\[
|\arg\lambda|>\theta
\]
if and only if
\[
c<H_\theta(a,b).
\]
Moreover,
\[
H_\theta(sa,s^2b)=s^3H_\theta(a,b).
\]
\end{lemma}

\begin{proof}
A positive-coefficient cubic has no nonnegative real root. A root can therefore leave the open angular region only through one of the rays \(\lambda=re^{\pm i\theta}\), \(r>0\). Substituting \(\lambda=re^{i\theta}\) and dividing the imaginary part by \(r\sin\theta\) gives
\[
Kr^2-2uar-b=0,
\]
whose unique positive solution is \(r_\theta(a,b)\). The real part then gives
\[
c=r^2(a+2ur)=H_\theta(a,b).
\]
Thus there is exactly one positive \(c\) at which a ray is touched. For \(c>0\) sufficiently small, one root is a small negative real root and the other two perturb the Hurwitz roots of \(\lambda^2+a\lambda+b\); for \(c\to\infty\), two roots approach the cube-root directions \(\pm\pi/3\), which violate \(\theta>\pi/3\). Hence stability is exactly \(c<H_\theta(a,b)\). Homogeneity follows by \(r_\theta(sa,s^2b)=sr_\theta(a,b)\).
\end{proof}

For \(2/3<\alpha<1\), put
\[
\theta=\frac{\alpha\pi}{2},
\qquad
u=\cos\theta,
\qquad
K=1-4u^2,
\]
and define
\[
r_\alpha(b)
=
\frac{u+\sqrt{u^2+Kb}}{K},
\qquad
h_\alpha(b)
=
r_\alpha(b)^2(1+2u\,r_\alpha(b)).
\]

\subsection{Exact elimination of the diagonal multiplier}

Define
\[
\boxed{
T_\alpha(\beta)
=
\min_{x\in\Delta_2^\circ}
\frac{h_\alpha(B_\beta(x))}
{x_1x_2x_3}.
}
\]
The minimum is attained in the interior because the numerator has a positive continuous extension to the closed simplex while \(x_1x_2x_3\to0\) on its boundary.

\begin{theorem}[Exact real-\(3\times3\) criterion]
\label{thm:C10}
Let \(-A\) be strict P and \(2/3<\alpha<1\). Then
\[
\boxed{
A\in\mathcal F_\alpha^{(3)}
\iff
\kappa<T_\alpha(\beta).
}
\]
\end{theorem}

\begin{proof}
For the simplex coordinate \(x\) associated with \(D\),
\[
\det(\lambda I-DA)
=
\lambda^3+s\lambda^2+s^2B_\beta(x)\lambda
+s^3\kappa x_1x_2x_3.
\]
Lemma~\ref{lem:cubic-boundary} and its homogeneity give
\[
\sigma(DA)\subset\Sigma_\alpha
\iff
\kappa
<
\frac{h_\alpha(B_\beta(x))}
{x_1x_2x_3}.
\]
Lemma~\ref{lem:simplex-reduction} identifies positive diagonal ratios with all \(x\in\Delta_2^\circ\), so the inequality holds for every \(D\succ0\) exactly when \(\kappa<T_\alpha(\beta)\).
\end{proof}

This is an explicit solution of the real three-dimensional Matignon instance of generalized positive-diagonal region stability. The abstract all-\(D\) framework itself is known \citep{Kushel2019,KushelPavani2022}.

\subsection{Interior classification}

\begin{lemma}[\(P_0\) necessity]
\label{lem:P0}
If \(A\in\mathcal F_\alpha^{(3)}\), then
\[
-a_{ii}\ge0,
\qquad
\det A[\{i,j\}]\ge0,
\qquad
-\det A>0.
\]
\end{lemma}

\begin{proof}
If \(a_{ii}>0\), keep the \(i\)-th row scaling fixed and send the other two to zero; the limiting matrix has a positive eigenvalue \(d_i a_{ii}\), which persists for nearby positive scalings. If an order-two principal minor is negative, send the remaining row scaling to zero; the limiting \(2\times2\) principal block has eigenvalues of opposite signs. Both contradict Matignon stability. Finally, at \(D=I\), the real eigenvalue of a real cubic must be negative, while the product of the remaining two eigenvalues is positive; hence \(\det A<0\).
\end{proof}

\begin{corollary}[Strict-P interior]
\label{cor:strictP-interior}
If
\[
A\in\operatorname{int}\mathcal F_\alpha^{(3)},
\]
then \(-A\) is strict P.
\end{corollary}

\begin{proof}
Lemma~\ref{lem:P0} gives nonnegative signed minors. A vanishing order-one minor can be perturbed immediately to the forbidden sign. A singular \(2\times2\) principal block can likewise be perturbed arbitrarily little to negative determinant: use a first-order adjugate perturbation when the block has rank one, and \(\operatorname{diag}(t,-t)\) when it is zero. Thus no such equality can occur at an interior point; the determinant inequality is already strict.
\end{proof}

Kellogg's theorem \citep{Kellogg1972} states that every eigenvalue \(\mu\) of a \(3\times3\) strict P-matrix satisfies
\[
|\arg\mu|<2\pi/3.
\]
Applied to \(-DA\), it yields
\[
|\arg\lambda(DA)|>\pi/3.
\]

\begin{theorem}[Interior of \(\mathcal F_\alpha^{(3)}\)]
\label{thm:interior-F}
For \(0<\alpha\le2/3\),
\[
\operatorname{int}\mathcal F_\alpha^{(3)}
=
\{A:-A\text{ strict P}\}.
\]
For \(2/3<\alpha<1\),
\[
\operatorname{int}\mathcal F_\alpha^{(3)}
=
\{A:-A\text{ strict P},\ \kappa<T_\alpha(\beta)\}.
\]
\end{theorem}

\begin{proof}
For \(\alpha\le2/3\), Kellogg gives sufficiency, Corollary~\ref{cor:strictP-interior} gives necessity, and strict P is open.

For \(\alpha>2/3\), necessity follows from Corollary~\ref{cor:strictP-interior} and Theorem~\ref{thm:C10}. Conversely, near any strict-P matrix the invariants vary continuously and remain in a compact positive box. Since
\[
F_{\alpha,\beta}(x)
\ge
\frac{h_\alpha(0)}{x_1x_2x_3},
\]
all nearby minimizers lie in one compact subset of \(\Delta_2^\circ\); hence \(T_\alpha(\beta)\) is locally continuous. The strict gap \(T_\alpha(\beta)-\kappa>0\) therefore persists under sufficiently small perturbations.
\end{proof}

\subsection{Cain endpoint and the exact fractional band}

As \(\alpha\to1^-\),
\[
h_\alpha(b)\to b,
\]
so
\[
T_\alpha(\beta)\to
T_1(\beta)
=
\min_{x\in\Delta_2^\circ}
\frac{B_\beta(x)}{x_1x_2x_3}.
\]
Since
\[
\frac{B_\beta(x)}{x_1x_2x_3}
=
\frac{\beta_{23}}{x_1}
+\frac{\beta_{13}}{x_2}
+\frac{\beta_{12}}{x_3},
\]
Cauchy--Schwarz gives
\[
\boxed{
T_1(\beta)
=
\left(
\sqrt{\beta_{12}}
+\sqrt{\beta_{13}}
+\sqrt{\beta_{23}}
\right)^2.
}
\]
This is exactly Cain's real-\(3\times3\) D-stability boundary \citep{Cain1976}.

\begin{lemma}[Strict fractional expansion of Cain's boundary]
\label{lem:Talpha-gt-T1}
For \(2/3<\alpha<1\),
\[
T_\alpha(\beta)>T_1(\beta)
\qquad
\forall\beta>0.
\]
\end{lemma}

\begin{proof}
For \(r=r_\alpha(b)\),
\[
b=Kr^2-2ur,
\]
hence
\[
h_\alpha(b)-b
=
4u^2r^2+2ur(r^2+1)>0.
\]
Apply this pointwise to the fractional objective and then minimize.
\end{proof}

\begin{corollary}[Exact genuinely fractional band]
\label{cor:exact-band}
For \(2/3<\alpha<1\),
\[
\boxed{
\operatorname{int}\mathcal P_\alpha^{(3)}
=
\{A:-A\text{ strict P},\
T_1(\beta)<\kappa<T_\alpha(\beta)\}.
}
\]
For \(0<\alpha\le2/3\),
\[
\boxed{
\operatorname{int}\mathcal P_\alpha^{(3)}
=
\{A:-A\text{ strict P},\ \kappa>T_1(\beta)\}.
}
\]
\end{corollary}
